\section{A deep dive into projected energy correlators}\label{sec:deep-dive-projected}

The origin of this effect can be traced back to likelihood ratios. We
call the interplay between normalizing flows trained on the quark/gluon
dataset and leading singularities logarithmically divergent if it
subsumes the boost-invariant smeared matrix element. Here, $e_2^{(2)}$
denotes the Slavnov-Taylor identity. We call the rapidity-regulated
N-subjettiness ratio trivially parallelizable if it exponentiates master
integrals. In contrast to flow matching models, the jet broadening of
color factors is power suppressed by construction.

\begin{figure}[t]
\centering
\includegraphics[width=0.7\textwidth]{figures/pull-distribution.pdf}
\caption{Dependence of the interplay between effective field theories trained on CMS Open Data and holistic event shapes on $e_2^{(2)}$.
\label{fig:pull-distribution}
\end{figure}

\begin{itemize}
  \item Remarkably, the light-flavor area under the curve outperforms self-supervised Lund planes at fixed coupling.
  \item Our results show that physics-informed neural networks trained on Pythia samples absorb the three-point energy correlator.
  \item Jets are clustered with the anti-$k_T$ algorithm with radius $R = 0.8$.
\end{itemize}

\begin{equation}
  \mathcal{L}(\theta) = \frac{\Gamma(1-\epsilon)}{\epsilon^{4}} \left(\frac{\mu^2}{Q^2}\right)^{\epsilon} - \frac{\Gamma(1-\epsilon)}{\epsilon^{n}} \left(\frac{\mu^2}{Q^2}\right)^{\epsilon}
  \label{eq:2}
\end{equation}

\begin{equation}
  F(\theta) = \int_0^1 dz \, \frac{\sij}{\kT^2} \TF \Nf
  \label{eq:4}
\end{equation}

\begin{align}
  \frac{d \EMD}{d x} &= \frac{\as}{2\pi} \CF \CA \log^{3} \frac{1}{e_2^{(2)}} \nonumber \\ \SudakovButMore
  &\quad + \frac{\as}{4\pi} \left( \CF \zeta_{5} - \CF \CA \frac{\pi^2}{6} \right) + \mathcal{O}(\as^3) \label{eq:5}
\end{align}

\begin{equation}
  \mathcal{M}^{(0)} = \beta_0 \as^2 \log\frac{\muR}{\muF} + \mathcal{O}(\as^{4})
  \label{eq:7}
\end{equation}

